\scanpage{127}
\mspage{466}
Pour des raisons qui restent encore mystérieuses pour moi, dans mon propre
cas l'histoire des relations (tant conscientes qu'inconscientes) entre le moi
(``le patron'')\inserted{,} et ``le mâle'' et ``le féminin'' en ma personne (aussi bien dans le
``patron'' lui-même que dans l'``ouvrier'', qui l'un et l'autre sont tributaires du
double aspect yin-yang de toutes choses) - cette histoire a été plus mouvementée
qu'à l'accoutumée. J'y distingue trois périodes. La dernière rejoint dans un cer-\newline
tain sens la première, qui s'étend sur les cinq premières années de mon enfance.
Cette troisième période, que je peux appeler celle de la\quad m a t u r i t é\quad,
peut être vue comme une sorte de ``retour'' à cette enfance, ou comme de progres-\newline
sives retrouvailles avec l' ``é t a t\quad d'\quad e n f a n c e'', avec l'harmonie
des épousailles sans histoires du ``yin'' et du ``yang'' en mon être. Ces retrou-\newline
vailles ont commencé au mois de juillet 1976, à l'âge de quarante huit ans -
l'année même où j'ai fait la découverte (trois mois plus tard) d'un pouvoir
jusque là ignoré en moi, le pouvoir de méditation (*).

Les valeurs dominantes dans la personne de chacun de mes parents, tant ma
mère que mon père, étaient des valeurs yang: volonté, intelligence (au sens: \unsure{pui|puis-}\newline
sance intellectuelle), contrôle de soi, ascendant sur autrui, intransigeance,
``Konsequenz'' (qui signifie, en allemand, cohérence extrême dans (ou avec) ses
options, idéologiques notamment), ``idéalisme'' au niveau politique comme prati-\newline
que\ldots Chez ma mère, cette valorisation a pris dès son jeune âge une force exa-\newline
cerbée, c'était le revers d'une véritable haine qu'elle avait développée vis à
vis de ``la femme'' en elle (et à partir de là, vis-à-vis du féminin en général:\inserted{)}.
Cette haine en elle a fini par prendre une véhémence et une force d'autant plus
destructice, qu'elle est restée entièrement occultéesa vie durant.\inserted{(}Moi-même ai
fini par découvrir ces choses il y a cinq ans seulement, trois ans après que la
méditation apparaisse dans ma vie.) Dans un tel contexte parental, c'est un
mystère (et pourtant un fait qui ne fait aucun doute pour moi) que j'aie pu
pleinement m'épanouir pendant les cinq premières années de mon enfance - \unsure{jusqu'a|jusqu'au}
moment de l'arrachement au milieu parental et de la destruction de ma famille
d'origine (formée de mes parents, de ma soeur plus agée, et de moi), de par la
volonté de ma mère et à la faveur (si on peut dire) des événements politiques de
l'année 1933.

\figure{\begin{tikzpicture}\draw (0,0)--(3.3,0);\end{tikzpicture}}
(*) Voir les deux sections ``Désir et méditation'' et ``L'émerveillement'', n°s 36
et 37.
\mspage{126}
